\section{Introduction}

Water pumps are vital in modern societies, supporting the supply of potable water, agricultural irrigation, and a wide range of industrial and urban activities. Despite their importance, pump malfunctions and leakages represent a significant challenge, as failures can lead to substantial economic and environmental impacts, such as water waste, service disruptions, and increased maintenance and repair costs \cite{che2021water, moleda2020predictive, jakobsen2023detecting}. Consequently, the early detection of pump failures is critical to mitigating these adverse effects, enabling timely maintenance actions, and ensuring the efficient, reliable, and sustainable operation of water pumping systems.

Conventional approaches to detect pump failure often focus on hydraulic parameters such as pressure and flow \cite{buffa2021advanced,jiang2020novel}. Although useful, these methods may not detect issues that primarily affect pump mechanical integrity, such as internal leaks or component failure \cite{barrientos2023water}. From an operational perspective, acoustic signal analysis provides a complementary and distinctive perspective on pump operation, as mechanical degradation typically manifests as subtle changes in the emitted sound pattern, enabling early detection of anomalies that may indicate impending problems.

Recent advancements in artificial intelligence have led to the development of several methods for detecting pump faults and anomalies using acoustic signals \cite{sun2023pump,pu2023research}. However, implementation in corporate and municipal water supply systems often depends on dedicated servers and continuous data transmission from field devices, which presents challenges when pumps and transducers are located in remote or inaccessible areas \cite{olsson2021urban}. To enable efficient failure detection in the field, TinyML offers a promising solution for acoustic signal analysis by enabling Machine Learning (ML) algorithms to run directly on low-power embedded devices, such as microcontrollers integrated into pumps \cite{abadade2023comprehensive}. By performing local, real-time signal processing without requiring cloud-based infrastructure, this approach reduces communication overhead, thereby enhancing system autonomy and enabling faster responses to emerging failures.

This paper proposes a TinyML-based fault-detection system for monitoring pump operation in water supply networks that distinguishes between normal and abnormal operating conditions using acoustic signals. The approach evaluates and compares various acoustic signal preprocessing techniques and integrates a deep learning model into a resource-constrained embedded platform, specifically the Arduino Nano 33 BLE Sense, to enable fully on-device, real-time inference. The main contributions include: (i) an end-to-end workflow for developing and deploying a TinyML model for acoustic-based pump fault detection, (ii) a comparative evaluation of two audio preprocessing techniques regarding memory footprint, latency, and inference time on an embedded device, and (iii) the design, implementation, and experimental validation of a fully embedded system for water pump fault identification.
